\begin{opinion}
    El trabajo \textbf{“DermaUH versión 3.0: Integración de nuevos modelos de apoyo al diagnóstico de 
    lesiones de cáncer de piel”} da respuesta a una de las tareas del proyecto “Desarrollo de 
    aplicaciones basadas en análisis imagenológico en apoyo al diagnóstico y tratamiento de 
    lesiones cutáneas”. Este proyecto está asociado al programa Nacional 6 “Telecomunicaciones
    e Informatización de la sociedad”.

    En este trabajo se rediseñan partes para corregir limitaciones de versiones previas, y 
    la introducción de nuevos modelos de diagnóstico desarrollados por el grupo. En un 
    inicio nos propusimos reentrenar esos modelos con imágenes de pacientes cubanas que se 
    están coleccionando con la aplicación móvil del Ecosistema. Las condiciones actuales de 
    trabajo de los especialistas a pesar de que la apk fue desarrollada teniendo en cuenta 
    entornos de no conectividad o conectividad inestable. Hasta hace pocos días la cantidad 
    de imágenes no es suficiente y contiene una mayoría considerable de carcinomas basales, 
    una de las lesiones que se consideran. Este objetivo tuvo que ser pospuesto y se 
    introdujeron los modelos entrenados con las imágenes disponibles. El texto analiza 
    herramientas comerciales existentes y repositorios académicos, concluyendo que estas 
    presentan sesgos hacia pieles claras o costos inaccesibles para el contexto de Cuba. 
    Mediante una metodología Scrum y una arquitectura de software organizada en capas, el 
    sistema busca optimizar el rendimiento y la seguridad en entornos de conectividad 
    limitada. Finalmente, se describen los requerimientos funcionales y el diseño de la 
    base de datos necesaria para gestionar pacientes, lesiones y diagnósticos dermatoscópicos
    de manera eficiente.

    La tesis de José Miguel es esencialmente un trabajo de ingeniería de software sin embargo,
    ha tenido que trabajar con los modelos de clasificación con diferentes arquitecturas, 
    clasificación basada en redes neuronales, arquitectura Transformer, clasificación basada 
    en extracción de características y un modelo explicativo. Además, tuvo que trabajar con 
    la Base de Imágenes de DermaUH. Todos los resultados obtenidos, incluso en estas 
    condiciones en que se ha trabajado han sido posibles gracias a su capacidad de asimilar 
    nuevos conocimientos, nuevas habilidades mostrando creatividad y capacidad de trabajo 
    en equipo consultando no solo con sus tutores sino con los miembros del equipo DermaUH.

    Consideramos que el principal valor de este trabajo radica en que aporta 
    significativamente a la evolución del ecosistema DermaUH para lograr una herramienta de 
    diagnóstico asistido, si bien aún perfectible, verdaderamente adaptada a la realidad 
    epidemiológica y tecnológica de Cuba. Algunos de los resultados más destacados incluyen:
    
    \begin{itemize}
        \item Optimización del rendimiento y seguridad: Se logró mejorar significativamente el servicio al usuario, reduciendo los tiempos de carga del frontend y garantizando una comunicación cifrada y segura entre cliente y servidor.
        \item DermaUH se concibe como una herramienta gratuita para los dermatólogos cubanos, eliminando las barreras económicas de suscripciones costosas (como las de DermEngine o VisualDx) con las que se compara en el trabajo y funcionando bajo condiciones de conectividad local.
    \end{itemize}

    José Miguel ha demostrado una notable madurez profesional y capacidad analítica durante todo el desarrollo de la tesis:
    \begin{itemize}
        \item Capacidad crítica: Realizó un análisis profundo del estado del arte, identificando con precisión las brechas de las herramientas comerciales actuales (costos prohibitivos y dependencia total de la nube) para proponer una solución viable en el contexto cubano.
        \item Autonomía en la toma de decisiones técnicas: El autor justificó de manera independiente la selección de una arquitectura en capas, priorizando el rendimiento y la mantenibilidad sobre patrones más complejos que podrían degradar la experiencia en condiciones de internet inestable.
        \item Creatividad en la solución de problemas: Ante las deficiencias detectadas en la versión 2.0, José Miguel no se limitó a realizar mejoras superficiales, sino que reestructuró el sistema utilizando la metodología ágil Scrum, dividiendo el trabajo en seis sprints estratégicos que abarcaron desde la gestión de roles hasta la integración de modelos de inferencia avanzados.
        \item Diseño integral: Su capacidad para modelar desde la base de datos (MER) hasta los protocolos de notificación por correo demuestra un dominio completo del ciclo de vida del desarrollo de software.
    \end{itemize}

    El trabajo de José Miguel Leyva De la Cruz no solo cumple con los objetivos académicos de 
    un ejercicio de conclusión de estudios, sino que ha demostrado una excepcional capacidad 
    para aplicar conocimientos avanzados de Ciencias de la Computación en la resolución de 
    problemas médicos críticos en nuestro país.

    \begin{itemize}
        \item \textbf{Lic. Alejandro Beltrán}
        \item \textbf{MSc Yoana Campbell}
        \item \textbf{DraC. Marta L. Baguer}
    \end{itemize}
\end{opinion}